% ===========================================================================
% OWNER: M3 (Zarpash Nasim); M1/M2 add their test rows by PR
% GRADED BY: Form 3 #8
% DRAW FROM: eval/tests, runner/tests, evidence/tests, CI run links
% Rules: mid-eval-report/CLAUDE.md. Do not edit a file you do not own.
% ===========================================================================
\chapter{Detailed Test Cases}\label{app:test-cases}

\section{Detailed Per-Case Results}\label{appc:sec:percase}
Table \ref{ch8:tab:fixture-results} in Chapter 7 gives the summary numbers for all seven cases.
This section adds the exact assertion each victim makes and the measured counts the summary
table does not have room for, all read directly from the committed reports in
\texttt{eval/reports/*.json} and the fixture source under \texttt{fixtures/od-fixture/src/}.

\begin{longtable}{p{1.6cm}p{5.4cm}p{6.9cm}}
\caption{What each victim actually asserts, and the real measured counts}\\
\toprule
\textbf{Case} & \textbf{Assertion} & \textbf{Measured} \\
\midrule
\endfirsthead
\toprule
\textbf{Case} & \textbf{Assertion} & \textbf{Measured} \\
\midrule
\endhead
F1 & \texttt{assertEquals(0, Config.mode)} --- fails when the polluter has set it to 1 & 20/20
  reproduced (Wilson 0.84--1.00), 0/20 alone; original order 3 tests, minimised to 2 \\
F2 & \texttt{assertFalse(FeatureFlags.isTurboEnabled())} --- the helper reads
  \texttt{odfixture.turbo} internally, not the test method itself & 20/20 reproduced (Wilson
  0.84--1.00), 0/20 alone; original order 6 tests, minimised to 2 \\
F3 & \texttt{assertFalse(Toggles.flagA \&\& Toggles.flagB)} --- true only once \emph{both}
  polluters have run & 20/20 reproduced (Wilson 0.84--1.00), 0/20 alone; original order 15
  tests, minimised to 3 (both polluters + victim) \\
F4 & \texttt{assertFalse(LateFlag.isSet)} --- the polluter class \texttt{ZzzLatePolluterTest}
  sorts after the victim alphabetically, so the natural order never triggers it & 20/20
  reproduced via shuffle seed 1 (Wilson 0.84--1.00), 0/20 alone; original order 3 tests,
  minimised to 2 \\
N1 & \texttt{assertEquals(2, 1)} --- unconditionally false & 20/20 alone-failures (deterministic
  by construction) \\
N2 & A random boolean (\texttt{new Random().nextBoolean()}) gates the assertion --- intermittent
  by design, independent of order & 9/20 (45\%) alone-failures in the committed report
  (\texttt{eval/reports/n2.json}); the exact count varies run to run by design. A separate,
  earlier verification run on Linux/JDK8 measured 10/20 (50\%) during the mechanism fix itself
  (Chapter 7's defects table) --- a different run, not this one, cited here only to say the fix
  is balanced on both platforms rather than one \\
N3 & Checks an environment variable (\texttt{FLAKETRACE\_N3\_NEVER\_SET}) nothing in this
  project's toolchain ever sets & 0/20 in the discovered order, 0/20 across 31 shuffled orders,
  0/20 alone --- never fails under FlakeTrace's control \\
\bottomrule
\end{longtable}

\section{CI Run Evidence}\label{appc:sec:ci}
Chapter 7 cites the real CI run confirming all tests pass: CI run \texttt{38155613996} (4 jobs,
\texttt{.github/workflows/ci.yml}) and Report run \texttt{38155613976} (1 job,
\texttt{.github/workflows/report.yml}), both on \texttt{main} commit \texttt{0fbad3e}. Table
\ref{appc:tab:ci} gives each job's real start/end time, fetched from the GitHub Actions API, not
estimated.

\begin{table}[H]\centering
\caption{Job timing for the cited CI run}\label{appc:tab:ci}
\begin{tabular}{p{2.6cm}p{2.2cm}p{2.2cm}p{2.0cm}}
\toprule
\textbf{Job} & \textbf{Started} & \textbf{Completed} & \textbf{Duration} \\
\midrule
\texttt{python-eval} & 16:28:38 & 16:28:48 & 10s \\
\texttt{fixture-build} & 16:28:38 & 16:29:19 & 41s \\
\texttt{evidence} & 16:28:38 & 16:29:34 & 56s \\
\texttt{runner} & 16:28:38 & 16:31:33 & 2m55s \\
\texttt{build-report} & 16:28:38 & 16:31:00 & 2m22s \\
\bottomrule
\end{tabular}
\end{table}

\noindent All five concluded \texttt{success} (UTC timestamps, 11 October 2026). The
\texttt{runner} job is the slowest, consistent with it being the only job running 144 tests
against real JVMs rather than a Python-only or static-analysis workload.

\section{Baseline Run Detail}\label{appc:sec:baseline}
Chapter 4's Acceptance Criteria (Table \ref{ch4:tab:acceptance}) and claim B11 cite the random-
order baseline's headline result; Table \ref{appc:tab:baseline} gives the real per-trial counts
behind it, from \texttt{eval/reports/baseline\_results.json}
(\texttt{eval/tools/run\_baseline\_real.py}, native JDK 24, 4 independent trials per case, base
seeds 0/100/200/300, \texttt{max\_runs=50}).

\begin{table}[H]\centering
\caption{Random-order baseline: shuffles needed to find the bug, per trial}\label{appc:tab:baseline}
\begin{tabular}{p{1.2cm}p{2.0cm}p{6.5cm}}
\toprule
\textbf{Case} & \textbf{Pool size} & \textbf{Shuffles to find it (seeds 0/100/200/300)} \\
\midrule
F1 & 6 (1 polluter) & 1, 1, 3, 2 \\
F2 & 6 (1 polluter) & 1, 1, 3, 2 \\
F3 & 7 (2 polluters, both required) & 1, 2, 3, 1 \\
F4 & 6 (1 polluter, order-needs-reshuffle) & 1, 1, 3, 2 \\
\bottomrule
\end{tabular}
\end{table}

\noindent Found the bug in all 16 trials (4 cases $\times$ 4 seeds). F1, F2 and F4's trials are
byte-for-byte identical because their candidate pools are the same shape (6 single-method
classes) --- the shuffle only depends on pool structure, not test identity, so identical shapes
reliably produce identical sequences for the same seed. F3's different shape (7 classes, two
required polluters) gives a different but similarly small result.
